% =====================================================================
%  CHƯƠNG 3. KIẾN TRÚC KỸ THUẬT VÀ CƠ CHẾ SUY LUẬN THEO IBM THINK
% =====================================================================
\chapter{Kiến trúc kỹ thuật và cơ chế suy luận theo IBM Think}

\section{Chu trình nhận thức ba giai đoạn (IBM Cognitive Cycle)}
IBM Think \cite{ibm_agents} mô tả hoạt động của một AI Agent qua ba giai đoạn khép kín (Hình~\ref{fig:ch3-cycle}).
\begin{enumerate}
  \item \textbf{Khởi tạo mục tiêu và lập kế hoạch (Goal Initialization \& Planning).} Con người hoặc hệ thống đặt mục tiêu
        và các ràng buộc; tác tử phân rã yêu cầu thành những tác vụ con có thể đo lường được. Với tác vụ đơn giản, bước lập
        kế hoạch có thể lược bỏ.
  \item \textbf{Suy luận với công cụ sẵn có (Reasoning with Available Tools).} Tác tử không biết hết mọi thứ, nên nó
        dùng công cụ (tìm kiếm web, API, cơ sở dữ liệu, thậm chí các tác tử khác) để bổ sung thông tin còn thiếu, rồi cập
        nhật kế hoạch dựa trên dữ liệu mới.
  \item \textbf{Học hỏi và phản chiếu (Learning \& Reflection).} Tác tử tiếp nhận phản hồi, từ người dùng hoặc từ tác tử
        khác, lưu những gì đã học vào bộ nhớ và tinh chỉnh kế hoạch theo vòng lặp (iterative refinement) để không lặp lại
        lỗi cũ.
\end{enumerate}

\hinhtikz{ch3_ibm_cycle}{Chu trình nhận thức ba giai đoạn của AI Agent theo IBM Think.}{ch3-cycle}

\section{Đối chiếu hai mô hình suy luận: ReAct và ReWOO}
IBM Think \cite{ibm_agents} chỉ ra hai trường phái suy luận then chốt định hình hiệu năng của tác tử
(Hình~\ref{fig:ch3-react-rewoo}).

\subsection{ReAct (Reasoning and Action)}
ReAct \cite{yao_react} yêu cầu mô hình xen kẽ \emph{suy nghĩ} và \emph{hành động}: sau mỗi hành động, tác tử quan sát kết
quả rồi mới suy nghĩ bước kế tiếp. Vòng lặp có dạng
\[
  \text{Thought} \;\rightarrow\; \text{Action} \;\rightarrow\; \text{Observation} \;\rightarrow\; \text{Thought}
  \;\rightarrow\; \cdots \;\rightarrow\; \text{Final Answer}.
\]
Ưu điểm là thích nghi linh hoạt: nếu một công cụ trả về lỗi, bước suy nghĩ tiếp theo sẽ thấy và đổi hướng. Nhược điểm là
mỗi bước quan sát đòi hỏi một lần gọi LLM với toàn bộ ngữ cảnh tích lũy, nên độ trễ và chi phí token tăng theo số bước.

\subsection{ReWOO (Reasoning Without Observation)}
ReWOO \cite{xu_rewoo} tách pha lập kế hoạch khỏi pha thực thi và gồm ba mô-đun:
\begin{itemize}
  \item \textbf{Planner} lập trước toàn bộ kế hoạch, kể cả các lời gọi công cụ, dùng biến giữ chỗ (\#E1, \#E2,\dots)
        cho kết quả chưa có.
  \item \textbf{Worker} thực thi các lời gọi công cụ, có thể song song khi không phụ thuộc nhau.
  \item \textbf{Solver} tổng hợp kế hoạch và các bằng chứng thu được thành câu trả lời cuối.
\end{itemize}
Vì LLM chỉ được gọi ở pha lập kế hoạch và pha tổng hợp, ReWOO tiết kiệm token và giảm độ trễ. Đổi lại, kế hoạch đã lập
không tự thay đổi khi kết quả trung gian bất ngờ.

\hinhtikz{ch3_react_rewoo}{So sánh cơ chế vận hành: ReAct (vòng lặp suy nghĩ -- hành động -- quan sát liên tục) và ReWOO
(lập kế hoạch trước, thực thi song song, tổng hợp sau).}{ch3-react-rewoo}

\begin{table}[!htbp]\centering
\caption{Đối chiếu ReAct và ReWOO.}\label{tab:ch3-react-rewoo}
\small
\begin{tabularx}{\linewidth}{L{3.6cm}XX}
\toprule
\textbf{Tiêu chí} & \textbf{ReAct} & \textbf{ReWOO} \\
\midrule
Thứ tự suy luận -- thực thi & Xen kẽ từng bước & Lập kế hoạch toàn bộ trước khi thực thi \\
Số lần gọi LLM & Một lần sau mỗi quan sát & Hai lần (Planner và Solver) \\
Khả năng thích nghi & Cao, đổi hướng ngay khi gặp kết quả bất ngờ & Thấp hơn, kế hoạch cố định \\
Chi phí token, độ trễ & Cao, tăng theo số bước & Thấp, công cụ có thể chạy song song \\
Phù hợp & Bài toán khám phá, môi trường biến động & Quy trình có cấu trúc, cần tiết kiệm chi phí \\
\bottomrule
\end{tabularx}
\end{table}

\section{Các mô hình suy luận nâng cao khác}
Ngoài ReAct và ReWOO, các tác tử hiện đại còn áp dụng thêm nhiều mẫu suy luận (Hình~\ref{fig:ch3-patterns}):
\begin{itemize}
  \item \textbf{Plan-and-Solve:} lập kế hoạch trừu tượng ở mức cao trước, sau đó chi tiết hóa và giải từng bước.
  \item \textbf{LLM Compiler:} biểu diễn các tác vụ thành đồ thị phụ thuộc và thực thi song song những nhánh độc lập,
        giống trình biên dịch tối ưu luồng thực thi.
  \item \textbf{Reflection:} một tác tử (hoặc chính tác tử đó ở vai trò khác) đóng vai phản biện độc lập, rà soát kết
        quả và yêu cầu sửa trước khi hoàn tất.
\end{itemize}

\hinhtikz{ch3_patterns}{Ba mẫu suy luận nâng cao: Plan-and-Solve, LLM Compiler (đồ thị tác vụ song song) và Reflection
(tác tử thực thi -- tác tử phản biện).}{ch3-patterns}
